\documentclass[11pt]{article}
\usepackage{amsmath,amssymb,amsthm}
\usepackage[margin=1in]{geometry}

\newtheorem{theorem}{Theorem}
\newtheorem{lemma}[theorem]{Lemma}

\title{Problem 438: Integer Part of Polynomial Equation's Solutions}
\author{Project Euler Solutions}
\date{}

\begin{document}
\maketitle

\section{Problem Statement}
For an $n$-tuple of integers $t = (a_1, \ldots, a_n)$, let $(x_1, \ldots, x_n)$ be the real solutions of $x^n + a_1 x^{n-1} + \cdots + a_n = 0$, sorted so that $\lfloor x_i \rfloor = i$ for $1 \leq i \leq n$. Define $S(t) = \sum |a_i|$. Find $\sum_t S(t)$ for $n = 7$.

\section{Mathematical Foundation}

\begin{theorem}[Vieta's Formulas]
If $x_1, \ldots, x_n$ are the roots of $x^n + a_1 x^{n-1} + \cdots + a_n = 0$, then $a_k = (-1)^k e_k(x_1, \ldots, x_n)$, where $e_k$ is the $k$-th elementary symmetric polynomial.
\end{theorem}

\begin{proof}
Expand $\prod_{i=1}^n (x - x_i)$ and match coefficients.
\end{proof}

\begin{lemma}[Fractional Part Decomposition]
Writing $x_i = i + f_i$ with $f_i \in [0, 1)$, the constraint $\lfloor x_i \rfloor = i$ is equivalent to $f_i \in [0, 1)$.
\end{lemma}

\begin{proof}
$\lfloor x_i \rfloor = i$ iff $i \leq x_i < i + 1$ iff $0 \leq f_i < 1$.
\end{proof}

\begin{theorem}[Integrality Constraints]
The coefficients are all integers iff $e_k(1+f_1, 2+f_2, \ldots, n+f_n) \in \mathbb{Z}$ for all $1 \leq k \leq n$.
\end{theorem}

\begin{proof}
Direct from Vieta's formulas with the substitution $x_i = i + f_i$.
\end{proof}

\begin{lemma}[Sum Constraint]
$\sum_{i=1}^n f_i \in \{0, 1, \ldots, n-1\}$.
\end{lemma}

\begin{proof}
$e_1 = n(n+1)/2 + \sum f_i \in \mathbb{Z}$ forces $\sum f_i \in \mathbb{Z}$. Since $0 \leq f_i < 1$, we have $0 \leq \sum f_i < n$.
\end{proof}

\begin{theorem}[Newton's Identity Reduction]
The $e_k$ are determined by the power sums $p_k = \sum x_i^k$ via:
\[
k \cdot e_k = \sum_{j=1}^{k} (-1)^{j-1} e_{k-j}\, p_j, \quad e_0 = 1.
\]
All $e_k \in \mathbb{Z}$ iff all $p_k \in \mathbb{Z}$.
\end{theorem}

\begin{proof}
Newton's identities follow from logarithmic differentiation of $\prod(1 - x_i t)$. The equivalence of integer $e_k$ and integer $p_k$ follows by induction.
\end{proof}

\begin{lemma}[Finite Search Space]
For fixed $n$, the number of valid fractional part tuples is finite, as the $n$ integrality conditions constrain $n$ unknowns in $[0,1)^n$ to a discrete set.
\end{lemma}

\begin{proof}
Each integrality condition reduces continuous degrees of freedom. The bounded domain $[0,1)^n$ and discrete target $\mathbb{Z}$ ensure finitely many solutions.
\end{proof}

\section{Algorithm}
\begin{enumerate}
    \item For each $k \in \{0, 1, \ldots, n-1\}$ (the integer value of $\sum f_i$):
    \item Recursively enumerate $f_1, \ldots, f_n \in [0, 1)$ satisfying $\sum f_i = k$ and all $e_j \in \mathbb{Z}$.
    \item Use Newton's identities for early pruning of partial assignments.
    \item For each valid tuple, compute $S(t) = \sum |a_k|$ via Vieta's formulas.
    \item Sum all $S(t)$.
\end{enumerate}

\section{Complexity Analysis}
\begin{itemize}
    \item \textbf{Time:} Bounded by the (small) number of valid tuples; each verified in $O(n^2)$.
    \item \textbf{Space:} $O(n)$ for the recursion stack.
\end{itemize}

\section{Answer}

\[
\boxed{2046409616809}
\]

\end{document}
